\section{Experiments}\label{sec:experiments}

We evaluate NELU as a drop-in GELU replacement, keeping all other hyperparameters identical.
For CNNs, the normalization axis is global ($\mathrm{dim}{=}(1{,}2{,}3)$); for Transformers, it is per-token ($\mathrm{dim}{=}{-}1$), matching the scope of the downstream mixing operation (Appendix~\ref{app:ablation}).
ReLU is included as a reference for CNNs; for Transformers and language models, where GELU is the standard, we compare GELU vs.\ NELU directly.

%% ============================================================
\subsection{CIFAR-10/100: CNN architectures}\label{sec:cifar}
%% ============================================================

We compare ReLU, GELU, and NELU on CIFAR-10/100~\citep{krizhevsky2009learning} across six CNN architectures: ResNet-20/56~\citep{he2016deep}, WRN-28-10~\citep{zagoruyko2016wide}, DenseNet-100~\citep{huang2017densely}, MobileNet-v2~\citep{sandler2018mobilenetv2}, and ShuffleNet-v1~\citep{zhang2018shufflenet} (3 seeds each).

\begin{table}[t]
\centering
\caption{\textbf{CIFAR-10 / CIFAR-100 test accuracy (\%).} Mean $\pm$ std, 3 seeds.}
\label{tab:cifar}
\vspace{0.5em}
\small
\begin{tabular}{lcccccc}
\toprule
 & \multicolumn{3}{c}{CIFAR-10} & \multicolumn{3}{c}{CIFAR-100} \\
\cmidrule(lr){2-4} \cmidrule(lr){5-7}
Architecture & ReLU & GELU & NELU & ReLU & GELU & NELU \\
\midrule
ResNet-20    & -- & -- & -- & -- & -- & -- \\
ResNet-56    & -- & -- & -- & -- & -- & -- \\
WRN-28-10    & -- & -- & -- & -- & -- & -- \\
DenseNet-100 & -- & -- & -- & -- & -- & -- \\
MobileNet-v2 & -- & -- & -- & -- & -- & -- \\
ShuffleNet-v1 & -- & -- & -- & -- & -- & -- \\
\bottomrule
\end{tabular}
\end{table}

%% ============================================================
\subsection{ImageNet: ViT-B/L}\label{sec:imagenet}
%% ============================================================

We train ViT-B and ViT-L~\citep{dosovitskiy2021image} on ImageNet-1k~\citep{russakovsky2015imagenet} following the AugReg recipe of \citet{steiner2022how}, swapping only the activation function.

\begin{table}[t]
\centering
\caption{\textbf{ImageNet top-1 accuracy (\%).}}
\label{tab:imagenet}
\vspace{0.5em}
\begin{tabular}{lccc}
\toprule
Model & Params & GELU & NELU \\
\midrule
ViT-B & 86M  & -- & -- \\
ViT-L & 304M & -- & -- \\
\bottomrule
\end{tabular}
\end{table}

%% ============================================================
\subsection{Language modeling: GPT-2}\label{sec:language}
%% ============================================================

To verify that the normalization principle transfers beyond vision, we train GPT-2 Small (124M), Medium (355M), and Large (774M)~\citep{radford2019language} from scratch on FineWeb-Edu 10B~\citep{penedo2024fineweb}, swapping GELU for NELU while keeping all other settings identical.

\begin{table}[t]
\centering
\caption{\textbf{Language modeling} (GPT-2, FineWeb-Edu 10B). Validation perplexity $\downarrow$.}
\label{tab:language}
\vspace{0.5em}
\begin{tabular}{lccc}
\toprule
Model & Params & GELU & NELU \\
\midrule
GPT-2 Small  & 124M & -- & -- \\
GPT-2 Medium & 355M & -- & -- \\
GPT-2 Large  & 774M & -- & -- \\
\bottomrule
\end{tabular}
\end{table}

%% ============================================================
\subsection{Corruption robustness}\label{sec:robustness}
%% ============================================================

Section~\ref{sec:consequences} predicts that when the test-time distribution of $\rho$ shifts, scale-dependent gates produce output directions that differ from those seen during training.
Our claim is not that scale-dependent gating is universally harmful, but that it becomes particularly fragile when pre-activation scale shifts at test time.
Image corruptions provide a natural test bed: they perturb feature statistics without changing the label space.
We evaluate on CIFAR-100-C and ImageNet-C~\citep{hendrycks2019robustness} (19 corruptions $\times$ 5 severities) without any retraining.
Table~\ref{tab:corruption} reports category-level means; per-corruption results are in Appendix~\ref{app:corruption_full}.

\begin{table}[t]
\centering
\caption{\textbf{Corruption robustness} (mean accuracy \% over 5 severities, grouped by category). \textbf{Bold} = best.}
\label{tab:corruption}
\vspace{0.5em}
\small
\begin{tabular}{lccc}
\toprule
 & \multicolumn{3}{c}{\textbf{CIFAR-100-C} (ResNet-20)} \\
\cmidrule{2-4}
Category & ReLU & GELU & NELU \\
\midrule
Noise (4)     & 21.94 & 21.06 & \textbf{24.37} \\
Blur (5)      & 41.40 & 41.89 & \textbf{44.36} \\
Weather (5)   & 51.25 & 51.99 & \textbf{52.97} \\
Digital (5)   & 46.29 & 46.76 & \textbf{48.85} \\
\midrule
\textbf{Overall (19)} & 41.18 & 41.44 & \textbf{43.60} \\
\bottomrule
\end{tabular}

\vspace{0.8em}

\begin{tabular}{lcc}
\toprule
 & \multicolumn{2}{c}{\textbf{ImageNet-C} (ViT-B)} \\
\cmidrule{2-3}
Category & GELU & NELU \\
\midrule
Noise (4)   & -- & -- \\
Blur (5)    & -- & -- \\
Weather (5) & -- & -- \\
Digital (5) & -- & -- \\
\midrule
\textbf{Overall (19)} & \textbf{--} & \textbf{--} \\
\bottomrule
\end{tabular}
\end{table}

NELU ranks first on all 19 CIFAR-100-C corruptions.
On noise corruptions---which most strongly shift the pre-activation scale distribution---GELU falls \emph{below} ReLU (21.06\% vs.\ 21.94\%), confirming that smooth gating's scale sensitivity can outweigh its benefits under distribution shift.
NELU recovers the robustness of scale invariance while retaining smooth gating (24.37\%).